% GENERATED FILE. Edit canonical lessons, then run npm run generate:books.
% canonical-source-hash: fnv1a64:2cbe9263f3d26dad
% canonical-lessons: IT-C466-afa, IT-C466-afoso, IT-C466-scottarsi, IT-C466-siccita, IT-C466-brezza

\chapter{Il caldo}
\label{ch:it-il-caldo}
\hlchaptermodality{466}

\begin{chapteropening}
\textbf{By the end of this chapter:} \emph{I can talk about hot weather: the sultry heat, muggy days, getting sunburnt, drought and a cool breeze.}

Complete the last lesson of chapter 466: I can talk about hot weather: the sultry heat, muggy days, getting sunburnt, drought and a cool breeze.
\end{chapteropening}

\section[l'afa]{l'afa — the sultry heat, the mugginess}

\label{lesson:IT-C466-afa}



\begin{quote}
Before the new one: say the Italian for the snail, then the Italian for to slither.
\end{quote}

\subsection*{You'll want to know: l'afa}
\textbf{l'afa} — \textquotedblleft{}the sultry heat, the mugginess\textquotedblright{}.

\begin{quote}
Ad agosto in città c'è un'afa terribile.
\end{quote}

\begin{quote}
\emph{In August the city is terribly muggy.}
\end{quote}

Heat with no wind and damp air.

\begin{cousinweb}
The origin of \emph{afa} is not settled; it may imitate the sound of gasping for breath.
\end{cousinweb}

\begin{grammarlens}[title={What you've built}]
One more word for telling, explaining and describing at B1.
\end{grammarlens}

\subsection*{Your turn}
\begin{itemize}
\raggedright
  \item \emph{Say it:} \emph{l'afa}
  \item \emph{Say it:} \emph{l'afa}, once more
  \item \emph{Say it:} a sentence of your own with \emph{l'afa}
  \item \emph{Recall:} say the Italian for the snail, then the Italian for to slither, then say \emph{l'afa} again
\end{itemize}

\begin{tcolorbox}[breakable,colback=teal!4,colframe=teal!35!black,title={Before you move on}]
What does l'afa mean? (\textquotedblleft{}The sultry heat, the mugginess\textquotedblright{}.) Say it once more.
\end{tcolorbox}
\section[afoso]{afoso — sultry, muggy}

\label{lesson:IT-C466-afoso}



\begin{quote}
Before the new one: say the Italian for to slither, then the Italian for the sultry heat.
\end{quote}

\subsection*{You'll want to know: afoso}
\textbf{afoso} — \textquotedblleft{}sultry, muggy\textquotedblright{}.

\begin{quote}
Oggi è una giornata afosa.
\end{quote}

\begin{quote}
\emph{Today is a muggy day.}
\end{quote}

\emph{Un'estate afosa} — a sticky summer.

\begin{cousinweb}
\emph{afoso} is built on \emph{afa} 'sultry heat' with the ending \emph{-oso} 'full of'. The origin of \emph{afa} is not settled.
\end{cousinweb}

\begin{grammarlens}[title={What you've built}]
One more word for telling, explaining and describing at B1.
\end{grammarlens}

\subsection*{Your turn}
\begin{itemize}
\raggedright
  \item \emph{Say it:} \emph{afoso}
  \item \emph{Say it:} \emph{afoso}, once more
  \item \emph{Say it:} a sentence of your own with \emph{afoso}
  \item \emph{Recall:} say the Italian for to slither, then the Italian for the sultry heat, then say \emph{afoso} again
\end{itemize}

\begin{tcolorbox}[breakable,colback=teal!4,colframe=teal!35!black,title={Before you move on}]
What does afoso mean? (\textquotedblleft{}Sultry, muggy\textquotedblright{}.) Say it once more.
\end{tcolorbox}
\section[scottarsi]{scottarsi — to get sunburnt; to burn oneself}

\label{lesson:IT-C466-scottarsi}



\begin{quote}
Before the new one: say the Italian for the sultry heat, then the Italian for sultry.
\end{quote}

\subsection*{You'll want to know: scottarsi}
\textbf{scottarsi} — \textquotedblleft{}to get sunburnt; to burn oneself\textquotedblright{}.

\begin{quote}
Mettiti la crema o ti scotti.
\end{quote}

\begin{quote}
\emph{Put on sun cream or you will get burnt.}
\end{quote}

Hot drinks too: \emph{mi sono scottato con il caffè}.

\begin{cousinweb}
\emph{scottarsi} is inherited from Vulgar Latin \emph{excoctāre} 'to scorch' (EtymDB), built on Latin \emph{coquere} 'to cook'.
\end{cousinweb}

\begin{grammarlens}[title={What you've built}]
One more word for telling, explaining and describing at B1.
\end{grammarlens}

\subsection*{Your turn}
\begin{itemize}
\raggedright
  \item \emph{Say it:} \emph{scottarsi}
  \item \emph{Say it:} \emph{scottarsi}, once more
  \item \emph{Say it:} a sentence of your own with \emph{scottarsi}
  \item \emph{Recall:} say the Italian for the sultry heat, then the Italian for sultry, then say \emph{scottarsi} again
\end{itemize}

\begin{tcolorbox}[breakable,colback=teal!4,colframe=teal!35!black,title={Before you move on}]
What does scottarsi mean? (\textquotedblleft{}To get sunburnt; to burn oneself\textquotedblright{}.) Say it once more.
\end{tcolorbox}
\section[la siccità]{la siccità — the drought}

\label{lesson:IT-C466-siccita}



\begin{quote}
Before the new one: say the Italian for sultry, then the Italian for to get sunburnt.
\end{quote}

\subsection*{You'll want to know: la siccità}
\textbf{la siccità} — \textquotedblleft{}the drought\textquotedblright{}.

\begin{quote}
Quest'estate c'è stata una lunga siccità.
\end{quote}

\begin{quote}
\emph{This summer there was a long drought.}
\end{quote}

\emph{Secco} — dry; \emph{la siccità} is a long spell with no rain.

\begin{cousinweb}
\emph{siccità} is a learned borrowing from Latin \emph{siccitās} 'dryness' (EtymDB), from \emph{siccus} 'dry'.
\end{cousinweb}

\begin{grammarlens}[title={What you've built}]
One more word for telling, explaining and describing at B1.
\end{grammarlens}

\subsection*{Your turn}
\begin{itemize}
\raggedright
  \item \emph{Say it:} \emph{la siccità}
  \item \emph{Say it:} \emph{la siccità}, once more
  \item \emph{Say it:} a sentence of your own with \emph{la siccità}
  \item \emph{Recall:} say the Italian for sultry, then the Italian for to get sunburnt, then say \emph{la siccità} again
\end{itemize}

\begin{tcolorbox}[breakable,colback=teal!4,colframe=teal!35!black,title={Before you move on}]
What does la siccità mean? (\textquotedblleft{}The drought\textquotedblright{}.) Say it once more.
\end{tcolorbox}
\section[la brezza]{la brezza — the breeze}

\label{lesson:IT-C466-brezza}



\begin{quote}
Before the new one: say the Italian for to get sunburnt, then the Italian for the drought.
\end{quote}

\subsection*{You'll want to know: la brezza}
\textbf{la brezza} — \textquotedblleft{}the breeze\textquotedblright{}.

\begin{quote}
La sera arriva una brezza dal mare.
\end{quote}

\begin{quote}
\emph{In the evening a breeze comes in from the sea.}
\end{quote}

A gentle wind; \emph{la brezza marina} — the sea breeze.

\begin{cousinweb}
The origin of \emph{brezza} is not settled; it may be linked to French \emph{brise}.
\end{cousinweb}

\begin{grammarlens}[title={What you've built}]
One more word for telling, explaining and describing at B1.
\end{grammarlens}

\subsection*{Your turn}
\begin{itemize}
\raggedright
  \item \emph{Say it:} \emph{la brezza}
  \item \emph{Say it:} \emph{la brezza}, once more
  \item \emph{Say it:} a sentence of your own with \emph{la brezza}
  \item \emph{Recall:} say the Italian for to get sunburnt, then the Italian for the drought, then say \emph{la brezza} again
\end{itemize}

\begin{tcolorbox}[breakable,colback=teal!4,colframe=teal!35!black,title={Before you move on}]
What does la brezza mean? (\textquotedblleft{}The breeze\textquotedblright{}.) Say it once more.
\end{tcolorbox}
